\subsection{里斯自同态}

设 E 是向量空间，u 是 E 的自同态。

E 的向量子空间的序列 \((\operatorname{Ker}(u^{n}))_{n\in\mathbf{N}}\) 是递增的；它们的并是 E 的一个在 u 下稳定的子空间，称为 u 的幂零空间。

E 的向量子空间的序列 \((\mathrm{Im}(u^{n}))_{n\in\mathbf{N}}\) 是递减的；它们的交是 E 的一个在 u 下稳定的子空间，称为 u 的余幂零空间。

引理 2. --- 设 E 是向量空间，u 是 E 的一个有指标的自同态（A, V, p. 126）。

如果下列性质中有两条满足，那么第三条也满足：

\begin{enumerate}
\def\labelenumi{(\roman{enumi})}
\item
  自同态 u 的指标为 0；
\item
  u 的幂零空间 N 是有限维的；
\item
  u 的余幂零空间 I 有有限余维。
\end{enumerate}

对每个整数 \(n \geqslant 0\)，自同态 \(u^{n}\) 有指标，等于 \(n \operatorname{ind}(u)\)（A, V, p. 127, 引理 2）。自然整数序列 \((\dim(\operatorname{Ker}(u^{n})))_{n \in \mathbf{N}}\) 是递增的；要使它是平稳的，必须且只需 u 的幂零空间是有限维的。自然整数序列 \((\operatorname{codim}_{\operatorname{E}}(\operatorname{Im}(u^{n})))_{n \in \mathbf{N}}\) 是递增的；要使它是平稳的，必须且只需 u 的余幂零空间是有限维的。假设条件 (i) 满足；关系式

\[
\dim (\operatorname{Ker} (u ^ {n})) - \operatorname{codim} _ {\mathrm{E}} (\operatorname{Im} (u ^ {n})) = \operatorname{ind} (u ^ {n}) = n \operatorname{ind} (u) = 0
\]

于是蕴含条件 (ii) 与 (iii) 等价。反之，如果条件 (ii) 与 (iii) 满足，这个公式蕴含序列 \((n\operatorname{ind}(u))_{n\in\mathbf{N}}\) 是平稳的，因此 u 的指标为零。

引理 3. --- 设 E 是向量空间，u 是 E 的一个有指标的自同态。假设 u 的指标为零。设 N 是 u 的幂零空间，I 是它的余幂零空间。

\begin{enumerate}
\def\labelenumi{\alph{enumi})}
\item
  这时空间 E 是子空间 N 与 I 的直和（维尔--菲廷分解，参见 A, VIII, § 2, 第 2 目, p. 25）。
\item
  自同态 u 通过过渡到子空间，诱导了 N 的一个幂零自同态与 I 的一个自同构 \(u_{I}\)。
\item
  设 v 表示 E 的这样一个自同态：它在 N 上为零且在 I 上与 \(u_{I}^{-1}\) 重合。我们有 \(u \circ v = v \circ u\)，且自同态 \(1_{E} - u \circ v\) 是以 N 为像、以 I 为核的投影。
\item
  E 的每个与 u 交换的自同态都稳定 N 与 I，且与 v 交换。
\end{enumerate}

自同态 \(u\) 满足前一个引理的各条性质。

设 \(n \in \mathbf{N}\) 使得 \(\operatorname{Ker}(u^{n}) = N\)。那么 \(\operatorname{Ker}(u^{m}) = \operatorname{Ker}(u^{n})\) 对每个整数 \(m \geqslant n\) 成立；由于 \(\operatorname{ind}(u^{m}) = m \operatorname{ind}(u) = 0\)，这就蕴含 \(\operatorname{Im}(u^{m}) = \operatorname{Im}(u^{n})\) 对 \(m \geqslant n\) 成立，因而 \(\operatorname{I} = \operatorname{Im}(u^{n})\)。由 A, VIII, p. 25, prop. 2 可得断言 a) 与 b) 满足。

对每个 \(x \in N\)，我们有 \(v(x) = 0\) 与 \(u(x) \in N\)，因而 \(u(v(x)) = 0\) 与 \(v(u(x)) = 0\)。对每个 \(x \in I\)，我们有 \(u(v(x)) = v(u(x)) = x\)。这就蕴含 \(u \circ v = v \circ u\)，且 \(1_{E} - u \circ v\) 是以 N 为像、以 I 为核的投影。

设 w 是 E 的与 u 交换的自同态。设 \(n \in \mathbf{N}\) 使得 \(N = \operatorname{Ker}(u^{n})\) 且 \(I = \operatorname{Im}(u^{n})\)。设 \(x \in N\)；我们有 \(u^{n}(w(x)) = w(u^{n}(x)) = 0\)，因而 \(w(x) \in N\)。设 \(x \in I\)；存在 \(y \in E\) 使得 \(x = u^{n}(y)\)；那么 \(w(x) = w(u^{n}(y)) = u^{n}(w(y)) \in I\)。这就证明了 w 稳定 N 与 I。

最后证明 \(w\) 与 \(v\) 可交换。对每个 \(x\)（在 \(N\) 中），我们有 \(v(w(x)) = 0 = w(v(x))\)。设 \(x \in I\)；\(v(w(x))\) 与 \(w(v(x))\) 这两个 \(E\) 的元素都属于 \(I\)，且它们在 \(u\) 下的像都等于 \(w(x)\)，因而它们相等。由线性性，这就蕴含 \(v\) 与 \(w\) 交换。

注记. --- 假设存在整数 \(n \geqslant 0\) 使得 \(\operatorname{Ker}(u^{n}) = \operatorname{Ker}(u^{n+1})\)（相应地，\(\operatorname{Im}(u^{n}) = \operatorname{Im}(u^{n+1})\)）；那么 \(\operatorname{Ker}(u^{m}) = \operatorname{Ker}(u^{n})\)（相应地，\(\operatorname{Im}(u^{m}) = \operatorname{Im}(u^{n})\)）对每个整数 \(m \geqslant n\) 成立。

假设 \(u\) 的指标为零。对每个整数 \(n \geqslant 0\)，我们有

\[
n \operatorname{ind} (u) = \operatorname{ind} \left(u ^ {n}\right) = \operatorname{codim} _ {\mathrm{E}} \left(\operatorname{Im} \left(u ^ {n}\right)\right) - \dim \left(\operatorname{Ker} \left(u ^ {n}\right)\right).
\]

因此条件 \(\operatorname{Ker}(u^n) = \operatorname{Ker}(u^{n+1})\) 与 \(\operatorname{Im}(u^n) = \operatorname{Im}(u^{n+1})\) 等价。当它们满足时，向量空间 E 是 \(\operatorname{Ker}(u^n)\) 与 \(\operatorname{Im}(u^n)\) 的直和（A, VIII, p. 25, prop. 2），空间 \(\operatorname{Ker}(u^n)\) 是有限维的，且 \(u\) 通过过渡到子空间诱导了 \(\operatorname{Im}(u^n)\) 的一个自同构。

定义 2. --- 设 E 是分离的局部凸空间。E 的里斯自同态指的是 E 的这样的弗雷德霍姆自同态：其幂零空间是有限维的且其余幂零空间有有限余维。

命题 5. --- 设 E 是分离的局部凸空间。E 的每个里斯自同态的指标都为零。

这由定义和引理 2 得出。

例子. --- 1) E 的每个自同构都是 E 的里斯自同态。如果 E 是有限维的，\(\mathcal{L}(\mathrm{E})\) 的每个元素都是 E 的里斯自同态。

\begin{enumerate}
\def\labelenumi{\arabic{enumi})}
\setcounter{enumi}{1}
\item
   如果 \((\mathrm{E}_{i})_{i\in\mathrm{I}}\) 是分离局部凸空间的有限族，且对 I 中所有 i，\(u_{i}\) 是 \(\mathcal{L}(\mathrm{E}_{i})\) 的元素，那么自同态 \(u=\bigoplus_{i\in\mathrm{I}}u_{i}\)（它定义在局部凸空间 \(\bigoplus_{i\in\mathrm{I}}E_{i}\) 上）是里斯自同态当且仅当对所有 i，\(u_{i}\) 是。
\item
   设 E 是 R 上的分离局部凸空间，并设 \(\mathrm{E}_{(\mathbf{C})}\) 是它的复化。设 u 是 E 的自同态。要使 u 是里斯自同态，必须且只需它的复化 \(u_{(\mathbf{C})}\) 是 \(\mathrm{E}_{(\mathbf{C})}\) 的里斯自同态。
\end{enumerate}

命题 6. --- 设 E 是分离的局部凸空间，u 是 E 的里斯自同态。用 N 表示 u 的幂零空间，用 I 表示它的余幂零空间。

\begin{enumerate}
\def\labelenumi{\alph{enumi})}
\item
  E 的子空间 N 与 I 是闭的且在 u 下稳定，而空间 E 是它们的拓扑直和。
\item
  自同态 u 通过限制定义了 I 的一个自同构 \(u_{I}\)；
\item
   向量空间 N 是有限维的，且 u 通过限制定义了 N 的一个幂零自同态 \(u_{N}\)；
\item
  设 v 是 \(\mathcal{L}(\mathrm{E})\) 中这样的元素：它在 N 上为零且在 I 上与 \(u_{\mathrm{I}}^{-1}\) 重合。它是 E 的里斯自同态，且是与 u 交换的 u 的拟逆。E 的自同态 \(1_{\mathrm{E}} - u \circ v\) 是以 N 为像、以 I 为核的投影。\(\mathcal{L}(\mathrm{E})\) 中每个与 u 交换的元素都稳定 N 与 I，且与 v 交换。
\end{enumerate}

设 n 是自然数，使得 \(\operatorname{Ker}(u^{n}) = \operatorname{N}\) 且 \(\operatorname{Im}(u^{n}) = \operatorname{I}\)。由于 E 是分离的且 u 是连续的，\(\operatorname{N} = \operatorname{Ker}(u^{n})\) 在 E 中是闭的。由于 u 是弗雷德霍姆自同态，\(u^{n}\) 也是，且 \(\operatorname{I} = \operatorname{Im}(u^{n})\) 是闭的（III, p. 42, prop. 2）。由引理 2，空间 E 是 N 与 I 的直和，断言 a) 于是由 EVT, I, p. 15 的 prop. 3 得出。

由于 I 是闭的、在 E 中有有限余维且在 u 下稳定，由 u 诱导的映射 \(u_{I}\colon I \to I\) 是弗雷德霍姆自同态（III, p. 44, prop. 3）；它还是双射，因而是 I 的自同构（III, p. 42, prop. 2）。这就证明了 b)，其余断言由引理 2 得出。

自同态 \(v\)（它在命题的断言 \(d\)）中定义）是 \(u\) 的一个拟逆，称为 \(u\) 的典范拟逆。

里斯自同态享有弗雷德霍姆映射的某些稳定性性质。

命题 7. — 设 E 是分离的局部凸空间，u 是 E 的自同态。

\begin{enumerate}
\def\labelenumi{\alph{enumi})}
\item
   设 \(E_{1}\) 是 E 的有限余维闭子空间，且在 u 下稳定。设 \(u_{1}\) 是 \(\mathcal{L}(E_{1})\) 中在 \(E_{1}\) 上与 u 重合的元素。那么 u 是里斯自同态当且仅当 \(u_{1}\) 是；
\item
   假设 u 是指标为 0 的弗雷德霍姆自同态。设 \(\widehat{\mathbf{E}}\) 是 E 的完备化，并设 \(\widehat{u} \in \mathcal{L}(\widehat{\mathbf{E}})\) 是 u 的连续延拓。那么 u 是里斯自同态当且仅当 \(\widehat{u}\) 是；
\item
 给 E 的对偶 E' 赋以有界收敛拓扑、紧收敛拓扑或弱拓扑。如果 u 是 E 的里斯自同态，那么 \(^{t}\) u 是 E' 的里斯自同态。
\end{enumerate}
a) 要使 \(u\) 是指标为零的弗雷德霍姆自同态，必须且只需 \(u_1\) 是（III, p. 44, prop. 3）。另一方面，对每个 \(n \in \mathbf{N}\)，我们有 \(\mathrm{Ker}(u_1^n) = \mathrm{E}_1 \cap \mathrm{Ker}(u^n)\)，由此


\[
\dim \operatorname{Ker} (u _ {1} ^ {n}) \leqslant \dim \operatorname{Ker} (u ^ {n}) \leqslant \dim \operatorname{Ker} (u _ {1} ^ {n}) + \operatorname{codim} _ {\mathrm{E}} (\mathrm{E} _ {1}),
\]

因此序列 \((\dim\operatorname{Ker}(u^{n}))_{n\in\mathbf{N}}\) 有界当且仅当序列 \((\dim\operatorname{Ker}(u_{1}^{n}))_{n\in\mathbf{N}}\) 有界。于是，u 是 E 的里斯自同态当且仅当 \(u_{1}\) 是 \(E_{1}\) 的里斯自同态（引理 2）。

\begin{enumerate}
\def\labelenumi{\alph{enumi})}
\setcounter{enumi}{1}
\item
  由 III, p. 44 的 prop. 4，对每个 \(n \in \mathbf{N}\)，映射 \(\widehat{u}^n\) 是 \(\widehat{\mathrm{E}}\) 的弗雷德霍姆自同态，满足 \(\operatorname{Ker}(\widehat{u}^n) = \operatorname{Ker}(u^n)\) 且 \(\operatorname{ind}(\widehat{u}^n) = \operatorname{ind}(u^n) = n\operatorname{ind}(u)\)。于是，\(u\) 是里斯自同态当且仅当 \(\widehat{u}\) 是里斯自同态。
\item
  转置 \({}^{t}u\) 是 E' 的指标为 0 的弗雷德霍姆自同态（III, p. 43, 第 3 目）。由于 \((^{t}u)^{n}\) 的核是 \(u^{n}\) 的像的正交补（EVT, IV, p. 27, prop. 2），序列 \((\mathrm{Ker}({}^{t}u)^{n})_{n}\) 是平稳的。这就证明了 \({}^{t}u\) 是 E' 的里斯自同态。
\end{enumerate}

命题 8. --- 设 E 是分离的局部凸空间，u 是 E 的自同态。下列条件等价：

\begin{enumerate}
\def\labelenumi{(\roman{enumi})}
\item
  u 是 E 的里斯自同态；
\item
  存在 u 的拟逆 v，它与 u 交换；
\item
  存在 E 的闭向量子空间 \(E_{1}\)，它有有限余维且在 u 下稳定，使得 u 诱导了 \(E_{1}\) 的一个自同构。
\end{enumerate}

蕴含关系 (i) \(\Longrightarrow\) (ii) 由 prop. 6, d) 得出。

假设 \(u\) 有一个拟逆 \(v\)，它与 \(u\) 交换。设 \(\mathrm{E}_1\) 是 \(1_{\mathrm{E}} - u \circ v\) 的核；它是 E 的闭向量子空间，有有限余维，且在 \(u\) 与 \(v\) 下稳定。由于 \(u(v(x)) = v(u(x)) = x\) 对一切 \(x \in \mathrm{E}_1\) 成立，映射 \(u\) 与 \(v\) 诱导了 \(\mathrm{E}_1\) 的互逆的自同构。因此 (ii) 蕴含 (iii)。

最后，蕴含关系 (iii) \(\Longrightarrow\) (i) 由 Prop. 7 的断言 \(a\)）以及自同构是里斯自同态这一事实得出。

命题 9. --- 设 E 是分离的局部凸空间，u、v 是 \(\mathcal{L}(\mathrm{E})\) 中可交换的元素。下列条件等价：

\begin{enumerate}
\def\labelenumi{(\roman{enumi})}
\item
  自同态 u 与 v 都是里斯自同态。
\item
  自同态 \(u \circ v\) 是里斯自同态。
\end{enumerate}

假设 u 与 v 是里斯自同态；用 \(u'\) 与 \(v'\) 表示它们的典范拟逆。E 的自同态 u、v、\(u'\) 与 \(v'\) 两两交换（prop. 6, d)）。自同态 \(v' \circ u'\) 是 \(u \circ v\) 的拟逆且与 \(u \circ v\) 交换，因而 \(u \circ v\) 是里斯自同态（prop. 8）。

反之，假设 \(u \circ v\) 是里斯自同态，并设 w 是它的典范拟逆。由于 u 与 \(u \circ v\) 交换，由 prop. 6, d)，自同态 u 与 w 交换。同样 v 与 w 交换。因此 E 的自同态 u、v 与 w 两两交换；由此得出 E 的自同态 \(v \circ w\) 是 u 的拟逆，而 \(w \circ u\) 是 v 的拟逆。由 prop. 8，u 与 v 是里斯自同态。

\pandocbounded{\includegraphics[keepaspectratio]{images/part001_6cebf9a3.jpg}}

设 u 是 E 的里斯自同态。如果 v 是 E 的自同构，那么 E 的自同态 \(u \circ v\) 不一定是 E 的里斯自同态，即使 \(u \circ v - v \circ u\) 有有限秩也是如此。如果 h 是 E 的有限秩自同态，自同态 \(u + h\) 也不一定是里斯自同态（参见 III, p. 120, 习题 3）。

命题 10. --- 设 E 与 F 是局部凸空间，\(p \in \mathcal{L}(\mathrm{E};\mathrm{F})\) 且 \(q \in \mathcal{L}(\mathrm{F};\mathrm{E})\)。置 \(u = 1_{\mathrm{E}} - q \circ p\) 与 \(v = 1_{\mathrm{F}} - p \circ q\)。

\begin{enumerate}
\def\labelenumi{\alph{enumi})}
\item
  设 \(n \in \mathbf{N}\)。映射 \(p\) 通过限制定义了从 \(\mathrm{Ker}(u^n)\) 到 \(\mathrm{Ker}(v^n)\) 上的拓扑向量空间同构，并通过过渡到商定义了从 Coker( \(u^n\) ) 到 Coker( \(v^n\) ) 上的拓扑向量空间同构；
\item
  如果 u 的像在 E 中是闭的，那么 v 的像在 F 中是闭的；
\item
   如果 \(\operatorname{Ker}(u)\) 在 E 中有一个拓扑补，那么 \(\operatorname{Ker}(v)\) 在 F 中也有一个。如果 \(\operatorname{Im}(u)\) 在 E 中有一个拓扑补，那么 \(\operatorname{Im}(v)\) 在 F 中也有一个；
\item
  如果 u 是严格态射，那么 v 是严格态射。
\item
  如果 u 是 E 的自同构，那么 v 是 F 的自同构。
\item
  如果 \(u\) 是 E 的弗雷德霍姆自同态，那么 \(v\) 是 F 的弗雷德霍姆自同态，且我们有 \(\operatorname{ind}(u) = \operatorname{ind}(v)\)；
\item
  假设 E 与 F 是分离的。如果 u 是 E 的里斯自同态，那么 v 是 F 的里斯自同态。
\end{enumerate}

我们首先注意公式

\[
q \circ v = u \circ q, \quad v \circ p = p \circ u.\tag{6}
\]

然后我们证明两个引理。

引理 4. --- 设 \(n \in \mathbf{N}\)。置

\[
q _ {n} = \sum_ {i = 1} ^ {n} (- 1) ^ {i - 1} \binom {n} {i} q \circ (p \circ q) ^ {i - 1}.
\]

我们有 \(u^{n}=1_{E}-q_{n}\circ p\) 与 \(v^{n}=1_{F}-p\circ q_{n}\)。

\((q\circ p)^{i} = q\circ (p\circ q)^{i - 1}\circ p\) 对每个整数 \(i\geqslant 1\) 成立。然后我们计算

\[
u ^ {n} = (1 _ {\mathrm{E}} - q \circ p) ^ {n} = 1 _ {\mathrm{E}} - \sum_ {i = 1} ^ {n} (- 1) ^ {i - 1} {\binom {n} {i}} (q \circ p) ^ {i} = 1 _ {\mathrm{E}} - q _ {n} \circ p
\]

\[
v ^ {n} = (1 _ {\mathrm{F}} - p \circ q) ^ {n} = 1 _ {\mathrm{F}} - \sum_ {i = 1} ^ {n} (- 1) ^ {i - 1} \binom {n} {i} (p \circ q) ^ {i} = 1 _ {\mathrm{F}} - p \circ q _ {n},
\]

由此得到结论。

引理 5. --- 设 \(u'\) 是 \(\mathcal{L}(\mathrm{E})\) 的一个元素。置 \(v' = 1_{\mathrm{F}} + p \circ u' \circ q\)。如果 \(u'\) 是 \(u\) 的逆，那么 \(v'\) 是 \(v\) 的逆。如果 \(u'\) 是 \(u\) 的拟逆，那么 \(v'\) 是 \(v\) 的拟逆。如果 \(u'\) 与 \(u\) 交换，那么 \(v'\) 与 \(v\) 交换。

利用公式 (6)，我们计算

\[
v - v ^ {\prime} \circ v = - p \circ u ^ {\prime} \circ q \circ v = - p \circ u ^ {\prime} \circ u \circ q
\]

\[
v - v \circ v ^ {\prime} = - v \circ p \circ u ^ {\prime} \circ q = - p \circ u \circ u ^ {\prime} \circ q,
\]

由此得

\[
1 _ {\mathrm{F}} - v ^ {\prime} \circ v = p \circ (1 _ {\mathrm{E}} - u ^ {\prime} \circ u) \circ q
\]

\[
1 _ {\mathrm{F}} - v \circ v ^ {\prime} = p \circ (1 _ {\mathrm{E}} - u \circ u ^ {\prime}) \circ q,
\]

这就证明了引理。

有了引理 4，只需证明断言 \(a\)）对 \(n = 1\) 成立。依据公式 (6)，映射 \(p\) 与 \(q\) 通过过渡到子空间定义了连续线性映射 \(p'\) ： \(\operatorname{Ker}(u) \to \operatorname{Ker}(v)\) 与 \(q'\) ： \(\operatorname{Ker}(v) \to \operatorname{Ker}(u)\)，并通过过渡到商定义了连续线性映射 \(p''\) ： \(\operatorname{Coker}(u) \to \operatorname{Coker}(v)\) 与 \(q''\) ： \(\operatorname{Coker}(v) \to \operatorname{Coker}(u)\)。

由于 \(1_{\mathrm{E}} - u = q \circ p\) 且 \(1_{\mathrm{F}} - v = p \circ q\)，映射 \(p'\) 与 \(q'\) 一方面、\(p''\) 与 \(q''\) 另一方面，是互逆的同构。这就证明了 \(a\)）。

要使一个连续线性映射的像是闭的，必须且只需这个映射的余核是分离空间。于是 b) 由 a) 得出。

用 \(i\colon \operatorname{Ker}(u)\to \mathrm{E}\) 与 \(j\colon \operatorname{Ker}(v)\to \mathrm{F}\) 表示典范单射。假设 \(\operatorname{Ker}(u)\) 在 E 中有一个拓扑补，并设 \(r\colon \mathrm{E}\to \operatorname{Ker}(u)\) 是 \(i\) 的一个连续线性收缩。置 \(r' = p'\circ r\circ q\)。我们有 \(r'\in \mathcal{L}(\mathrm{F};\operatorname {Ker}(v))\) 且

\[
r ^ {\prime} \circ j = p ^ {\prime} \circ r \circ q \circ j = p ^ {\prime} \circ r \circ i \circ q ^ {\prime} = p ^ {\prime} \circ q ^ {\prime} = 1 _ {\operatorname{Ker} (v)},
\]

因而 \(r'\) 是 j 的一个连续线性收缩，且 \(\operatorname{Ker}(v)\) 在 F 中有一个拓扑补。这就证明了 c) 的第一个断言。第二个断言类似地证明，注意如果 \(s: \operatorname{Coker}(u) \to \operatorname{E}\) 是从 E 到 \(\operatorname{Coker}(u)\) 上的典范满射的一个连续线性截面，那么 \(p \circ s \circ q''\) 是从 F 到 \(\operatorname{Coker}(v)\) 上的典范满射的一个连续线性截面。

用 \(\overline{u}\colon\mathrm{E}/\mathrm{Ker}(u)\to\mathrm{Im}(u)\) 与 \(\overline{v}\colon\mathrm{F}/\mathrm{Ker}(v)\to\mathrm{Im}(v)\) 表示由 u 与 v 得到的连续线性双射。由公式 (6)，映射 p 通过过渡到商定义了一个连续线性映射 \(p_{1}\colon\mathrm{E}/\mathrm{Ker}(u)\to\mathrm{F}/\mathrm{Ker}(v)\)，且 p、q 通过过渡到子空间定义了连续线性映射 \(p_{0}\colon\mathrm{Im}(u)\to\mathrm{Im}(v)\) 与 \(q_{0}\colon\mathrm{Im}(v)\to\mathrm{Im}(u)\)。

用 \(t \colon \operatorname{Im}(v) \to \operatorname{F} / \operatorname{Ker}(v)\) 表示典范单射 \(\operatorname{Im}(v) \to \operatorname{F}\) 与典范满射 \(\operatorname{F} \to \operatorname{F} / \operatorname{Ker}(v)\) 的复合映射；它是连续的。我们有 \(\overline{v} \circ t = 1_{\operatorname{Im}(v)} - p_0 \circ q_0\) 与 \(\overline{v} \circ p_1 = p_0 \circ \overline{u}\)，由此

\[
\overline {{v}} \circ (t + p _ {1} \circ \overline {{u}} ^ {- 1} \circ q _ {0}) = \overline {{v}} \circ t + p _ {0} \circ q _ {0} = 1 _ {\mathrm{Im} (v)}.
\]

这证明了 \(t + p_{1} \circ \overline{u}^{-1} \circ q_{0}\) 是 \(\overline{v}\) 的逆双射。如果 u 是严格态射，那么映射 \(\overline{u}^{-1}\) 是连续的，\(\overline{v}^{-1}\) 亦然，这就证明 v 是严格态射。这就证明了 d)。

由引理 5，如果 \(u\) 是 E 的自同构，那么 \(v\) 是 F 的自同构。类似地，如果 \(u\) 是 E 的弗雷德霍姆自同态，那么 \(v\) 是 F 的弗雷德霍姆自同态，且这时有 \(\operatorname{ind}(u) = \operatorname{ind}(v)\)，依据 \(a\)）。这就证明了 \(e\)）与 \(f\)）。最后，如果 E 与 F 是分离的且 \(u\) 是 E 的里斯自同态，那么由引理 5 与 Prop. 8 的条件 (ii)，\(v\) 是 F 的里斯自同态；这就是断言 \(g\)）。

